\documentclass{article}
\usepackage{amsmath,amssymb,xurl}
\usepackage[hidelinks]{hyperref}
% Shared commands for notes in docs/paper-gaps/.

\DeclareUnicodeCharacter{2080}{\ifmmode{}_{0}\else\textsubscript{0}\fi}
\DeclareUnicodeCharacter{2081}{\ifmmode{}_{1}\else\textsubscript{1}\fi}
\DeclareUnicodeCharacter{2082}{\ifmmode{}_{2}\else\textsubscript{2}\fi}
\DeclareUnicodeCharacter{2099}{\ifmmode{}_{n}\else\textsubscript{n}\fi}

\newcommand{\Lean}{\textsc{Lean}}
\ExplSyntaxOn
\NewDocumentCommand{\leanid}{m}{
  \ifmmode
    \textsf{\detokenize{#1}}
  \else
    \group_begin:
    \urlstyle{sf}
    \tl_set:Nn \l_tmpa_tl {#1}
    \tl_replace_all:Nnn \l_tmpa_tl {\_} {_}
    \exp_args:NV \path \l_tmpa_tl
    \group_end:
  \fi
}
\ExplSyntaxOff
\providecommand{\tr}{\operatorname{tr}}
\newcommand{\tnleanIssueBase}{https://github.com/LionSR/TNLean/issues/}
\newcommand{\tnleanPrBase}{https://github.com/LionSR/TNLean/pull/}
\newcommand{\tnleanDocBase}{https://sirui-lu.com/TNLean/docs/}
\newcommand{\ghissue}[1]{\href{\tnleanIssueBase#1}{GitHub issue \##1}}
\newcommand{\ghpr}[1]{\href{\tnleanPrBase#1}{PR \##1}}
\newcommand{\doclink}[1]{\href{\tnleanDocBase#1.html}{\path{#1}}}

% Dirac notation and the identity operator, as in blueprint/src/macros/common.tex.
% \providecommand keeps note-local definitions authoritative.
\usepackage{dsfont}
\providecommand{\ket}[1]{|#1\rangle}
\providecommand{\bra}[1]{\langle #1|}
\providecommand{\braket}[2]{\langle #1 | #2 \rangle}
\providecommand{\ketbra}[2]{|#1\rangle\!\langle #2|}
\providecommand{\Id}{\mathds{1}}
\providecommand{\one}{\mathds{1}}
\providecommand{\C}{\mathbb{C}}
\providecommand{\MN}[1]{M_{#1}(\C)}
\providecommand{\MD}{M_D(\C)}

% Verdict marker: \gapnote{<kind>}{<status>}. Kind is the policy
% classification (clarification, local-correction, scope-restriction,
% unfaithful, false-source, open-gap); status is open, wip (elimination
% underway), resolved, or historical. Machine-read by the site index;
% prints a verdict line.
\newcommand{\gapnote}[2]{\par\noindent\textbf{Verdict:} #1 (#2).\par}

\title{Uniform Displacement Completeness and Missing Square Supports}
\author{}
\date{}
\begin{document}
\maketitle
\gapnote{scope-restriction}{open}

Schuch, Cirac, and P\'erez-Garc\'ia,
\href{https://arxiv.org/abs/1001.3807v3}{arXiv:1001.3807v3},
Lemma 6.14, source \path{Papers/1001.3807/paper_v3.tex}, lines 2199--2214,
uses endpoint-preserving string deformations. The full-support generated
homotopies of labelled torus paths are classified by lifted displacement.
The same uniform classification cannot be retained after excluding any
primal square-move support vertex.

The theorem
\leanid{TNLean.PEPS.torusDualHomotopy_displacement_complete_iff_support_univ}
proves that displacement classifies every same-endpoint pair of paths under
supported homotopy exactly when the allowed support set is the whole torus.
For each missing vertex $c$, a nonabelian seam word is based at $c$.
Allowed square moves preserve it, but the square surrounding $c$ has the
nontrivial commutator of two explicitly constructed transpositions on three
elements. Its displacement is zero, as is the constant path's displacement.
The two paths therefore cannot be related by the allowed moves.
The obstruction theorem is
\leanid{TNLean.PEPS.torusDualSquareLoop_not_homotopic_of_support_excludes}.

This statement concerns a uniform classification over all finite labelled
paths, including at torus periods zero, one and two. It does not claim that
individual equal-displacement paths fail to be homotopic, or give a finite
decision procedure for restricted homotopy. The excluded vertices are primal
supports of square moves; they are not a separately imposed restriction on
the visited dual vertices. The seam-word test group is independent of any
physical flux insertion group.

This nonabelian homotopy obstruction is distinct from an annular physical
Wilson-loop measurement. It proves no density inequality by itself and
does not equate physical states in different winding classes. The positive
restricted classification by all locally flat transports remains a separate
result; reducing it to a finite geometric normal form remains open.
The broader source interpretation stays under \ghissue{8703}, and
\ghissue{8947} is unchanged.
\end{document}
